%%%%%%%%%%%%%%%%%%%%%%%%%%%%%%%%%%%%%%%%%%%%%%%%%%%%%%%%%%%%%%%%%%%%%%%%%%%%%%%%%%
\begin{frame}[fragile]\frametitle{}
\begin{center}
{\Large Advanced Topics}

\vspace{0.3cm}
{\small For those who want to go deeper}
\end{center}
\end{frame}

%%%%%%%%%%%%%%%%%%%%%%%%%%%%%%%%%%%%%%%%%%%%%%%%%%%%%%%%%%%%%%%%%%%%%%%%%%%%%%%%%%
% -----------------------------------------------------------------------
% PLACEHOLDER: Topological Deep Learning
% The user will add a dedicated TDL section here.
% -----------------------------------------------------------------------
\begin{frame}[fragile]\frametitle{Topological Deep Learning --- Coming Soon}
  \begin{block}{Upcoming Section}
    A dedicated section on \textbf{Topological Deep Learning (TDL)} will be added here, covering:
    \begin{itemize}
      \item Topological Data Analysis (TDA) foundations: simplicial complexes, persistent homology
      \item Higher-order domains: simplicial, cellular, and hypergraph complexes
      \item Message passing on higher-order domains (TopoX, CIN, SCA)
      \item Connections between TDL and classical GNNs
      \item Applications: drug discovery, brain networks, material science
    \end{itemize}
  \end{block}
  \vspace{0.3cm}
  {\em Watch this space!}
\end{frame}

%%%%%%%%%%%%%%%%%%%%%%%%%%%%%%%%%%%%%%%%%%%%%%%%%%%%%%%%%%%%%%%%%%%%%%%%%%%%%%%%%%
\begin{frame}[fragile]\frametitle{}
\begin{center}
{\Large Lie Groups and Manifolds}
\end{center}
\end{frame}

%%%%%%%%%%%%%%%%%%%%%%%%%%%%%%%%%%%%%%%%%%%%%%%%%%%%%%%%%%%
\begin{frame}[fragile]\frametitle{Lie Groups: A Primer}
  \begin{itemize}
    \item A \textbf{Lie group} is both a \emph{group} (closed under composition, has identity and inverses) and a \emph{smooth manifold} (differentiable structure).
    \item This combination allows us to do calculus on groups --- compute gradients, define flows.
    \item Key to modelling \textbf{continuous symmetries} in data and physics.
    \item Examples: SO(3) (3D rotations), SE(3) (3D rigid motions), $\mathbb{R}^n$ (translations).
  \end{itemize}
\end{frame}

%%%%%%%%%%%%%%%%%%%%%%%%%%%%%%%%%%%%%%%%%%%%%%%%%%%%%%%%%%%
\begin{frame}[fragile]\frametitle{Core Properties of Lie Groups}
  \begin{itemize}
    \item \textbf{Differentiable structure}: enables calculus on groups (gradient descent on manifold).
    \item \textbf{Group properties}: closure, identity, inverse, associativity.
    \item \textbf{Lie algebra}: encodes local, infinitesimal transformations (the tangent space at the identity).
    \item \textbf{Exponential map}: takes a tangent vector (Lie algebra element) to a group element --- bridges local and global.
  \end{itemize}
\end{frame}

%%%%%%%%%%%%%%%%%%%%%%%%%%%%%%%%%%%%%%%%%%%%%%%%%%%%%%%%%%%
\begin{frame}[fragile]\frametitle{Lie Groups in Machine Learning}
  \begin{itemize}
    \item Build models that are \textbf{equivariant to continuous symmetry groups}.
    \item Image processing: invariance to rotation, translation, scale.
    \item \textbf{SE(3)-equivariant networks}: used for molecular geometry (atoms are points in 3D; the network must be invariant to how you rotate/translate the molecule).
    \item Applied in protein structure prediction (AlphaFold), materials science, robotics.
  \end{itemize}
\end{frame}

%%%%%%%%%%%%%%%%%%%%%%%%%%%%%%%%%%%%%%%%%%%%%%%%%%%%%%%%%%%
\begin{frame}[fragile]\frametitle{Lie Groups in Optimisation}
  \begin{itemize}
    \item Model parameters can evolve on curved manifolds rather than flat space.
    \item \textbf{Riemannian optimisation}: gradient descent on a Riemannian manifold.
    \item Lie algebra enables structured gradient computation on the group.
    \item \textbf{Exponential / log maps} link the tangent space (linear) to the group (curved).
    \item Used in: covariance matrix optimisation, SPD networks, normalising flows.
  \end{itemize}
\end{frame}

%%%%%%%%%%%%%%%%%%%%%%%%%%%%%%%%%%%%%%%%%%%%%%%%%%%%%%%%%%%
\begin{frame}[fragile]\frametitle{Differential \& Riemannian Manifolds}
  \begin{itemize}
    \item \textbf{Smooth manifold}: locally homeomorphic to $\mathbb{R}^n$; calculus is possible via local coordinate charts.
    \item \textbf{Riemannian manifold}: a smooth manifold with a \emph{metric tensor} that defines distances and angles.
    \item Enables geometry-aware learning: compute distances, geodesics, curvature.
    \item Basis for Riemannian batch normalisation, SPD networks, hyperbolic embedding.
  \end{itemize}
\end{frame}

%%%%%%%%%%%%%%%%%%%%%%%%%%%%%%%%%%%%%%%%%%%%%%%%%%%%%%%%%%%
\begin{frame}[fragile]\frametitle{Tangent Spaces and Geodesics}
  \begin{itemize}
    \item \textbf{Tangent space at a point}: local linear approximation of the manifold --- like zooming in until the surface looks flat.
    \item \textbf{Tangent vectors} act as directional derivatives on the manifold.
    \item \textbf{Geodesics}: shortest paths on a curved surface (analogue of straight lines in Euclidean space).
    \item Example: the shortest flight path between two cities follows a great circle on the sphere, not a straight line on a flat map.
  \end{itemize}
\end{frame}

%%%%%%%%%%%%%%%%%%%%%%%%%%%%%%%%%%%%%%%%%%%%%%%%%%%%%%%%%%%
\begin{frame}[fragile]\frametitle{Geomstats: Computing on Manifolds}
  \begin{itemize}
    \item Open-source Python library for computation on nonlinear manifolds.
    \item Supports Lie groups, Riemannian manifolds, and symmetric spaces.
    \item Follows the Scikit-Learn API for easy integration with standard ML pipelines.
    \item Ideal for hands-on exploration of geometric learning concepts.
  \end{itemize}
\end{frame}

%%%%%%%%%%%%%%%%%%%%%%%%%%%%%%%%%%%%%%%%%%%%%%%%%%%%%%%%%%%%%%%%%%%%%%%%%%%%%%%%%%
\begin{frame}[fragile]\frametitle{}
\begin{center}
{\Large Deeper into the 5G Framework}
\end{center}
\end{frame}

%%%%%%%%%%%%%%%%%%%%%%%%%%%%%%%%%%%%%%%%%%%%%%%%%%%%%%%%%%%%%%%%%%%%%%%%%%%%%%%%%%
\begin{frame}[fragile]\frametitle{Group Domain: Beyond Translation}

\begin{center}
\textbf{Homogeneous Spaces with Global Symmetries}
\end{center}

\begin{columns}
\begin{column}{0.5\textwidth}
\textbf{Examples:}
\begin{itemize}
\item \textbf{Sphere}: SO(3) rotations
\item \textbf{Scale space}: scaling transformations
\item \textbf{Roto-translation}: SE(2) group
\end{itemize}

\textbf{Group Convolution:}
$$[f *_G \psi](g) = \int_G f(h)\, \psi(g^{-1}h)\, dh$$

\textbf{Applications:}
\begin{itemize}
\item Spherical CNNs for 360\textdegree{} images
\item Scale-equivariant networks
\item Rotation-invariant 3D vision
\end{itemize}
\end{column}
\begin{column}{0.5\textwidth}
\begin{center}
\textit{[image: spherical CNN --- convolution kernel sliding over a globe]}

\vspace{0.3cm}
{\small Spherical convolution on global data}
\end{center}
\end{column}
\end{columns}

\begin{block}{Key Insight}
Group equivariance generalises translation equivariance to any continuous symmetry group.
\end{block}

\end{frame}

%%%%%%%%%%%%%%%%%%%%%%%%%%%%%%%%%%%%%%%%%%%%%%%%%%%%%%%%%%%%%%%%%%%%%%%%%%%%%%%%%%
\begin{frame}[fragile]\frametitle{Geodesic Domain: Manifold Learning}

\begin{center}
\textbf{Continuous Curved Spaces}
\end{center}

\begin{columns}
\begin{column}{0.45\textwidth}
\textbf{Key concepts:}
\begin{itemize}
\item Riemannian manifolds
\item Geodesic distances
\item Intrinsic geometry
\item Curvature effects
\end{itemize}

\textbf{Applications:}
\begin{itemize}
\item 3D shape analysis
\item Molecular surfaces
\item Brain cortex modelling
\end{itemize}
\end{column}
\begin{column}{0.55\textwidth}
\textbf{Geodesic Convolution:}
$$[f * \psi]_{\mathcal{M}}(x) = \int_{\mathcal{M}} f(y)\, \psi(d_{\mathcal{M}}(x,y))\, d\mu(y)$$

\begin{center}
\textit{[image: 3D triangular mesh with geodesic path highlighted]}
\end{center}
\end{column}
\end{columns}

\begin{block}{Diffeomorphism Invariance}
Functions invariant to smooth, invertible deformations of the manifold.
\end{block}

\end{frame}

%%%%%%%%%%%%%%%%%%%%%%%%%%%%%%%%%%%%%%%%%%%%%%%%%%%%%%%%%%%%%%%%%%%%%%%%%%%%%%%%%%
\begin{frame}[fragile]\frametitle{Gauge Domain: Vector Bundles}

\begin{center}
\textbf{The Most General Geometric Structure}
\end{center}

\textbf{Key concepts:}
\begin{itemize}
\item \textbf{Vector bundle}: at every point of a manifold, attach a vector space (the ``fibre'').
\item \textbf{Gauge transformation}: change the local coordinate frame for the fibre --- the physics should not depend on this choice.
\item \textbf{Parallel transport}: consistently ``carry'' a vector along a path on the manifold.
\end{itemize}

\textbf{Gauge Equivariant Convolution:}
$$[f *_{\text{gauge}} \psi](x) = \int_{\mathcal{M}} \Phi(x,y)\, f(y)\, \psi(x,y)\, d\mu(y)$$

where $\Phi(x,y)$ is the parallel transport operator.

\textbf{Applications:} Weather prediction (wind vectors on the sphere), fluid dynamics, electromagnetic fields, general relativity.

\begin{alertblock}{Ultimate Generalisation}
Gauge networks are the most general class of equivariant networks in the 5G framework.
\end{alertblock}

\end{frame}

%%%%%%%%%%%%%%%%%%%%%%%%%%%%%%%%%%%%%%%%%%%%%%%%%%%%%%%%%%%%%%%%%%%%%%%%%%%%%%%%%%
\begin{frame}[fragile]\frametitle{Current Challenges and Research Frontiers}

\begin{columns}
\begin{column}{0.5\textwidth}
\textbf{Theoretical:}
\begin{itemize}
\item Expressivity vs.\ efficiency trade-offs
\item Universal approximation on manifolds
\item Generalisation bounds for geometric networks
\end{itemize}

\textbf{Computational:}
\begin{itemize}
\item Scalability to very large graphs
\item GPU-efficient implementations
\item Distributed training strategies
\end{itemize}
\end{column}
\begin{column}{0.5\textwidth}
\textbf{Application:}
\begin{itemize}
\item Heterogeneous graph integration
\item Dynamic and temporal graphs
\item Uncertainty quantification
\end{itemize}

\textbf{Future directions:}
\begin{itemize}
\item Automated geometric architecture search
\item Continuous--discrete bridges
\item Quantum geometric learning
\item Causal geometric models
\end{itemize}
\end{column}
\end{columns}

\begin{alertblock}{The Grand Challenge}
\begin{center}
Can we develop a unified theory that \emph{automatically discovers} the right geometric structure for any data and task?
\end{center}
\end{alertblock}

\end{frame}
